% ============================================================
%  Capítulo 8 — Conclusão
% ============================================================
\chapter{CONCLUSÃO}
\label{cap:conclusao}

\section{Síntese das contribuições}
\label{sec:conclusao-contrib}

Este trabalho partiu de uma limitação concreta do estado da arte: o ByteDefender
\cite{nikkhahbahrami2025} demonstrou que bytecodes V8 são suficientemente
discriminativos para detectar funções de browser fingerprinting com
F1~$\approx$~0,95, mas o modelo é restrito a um único motor JavaScript e não
transfere para o Firefox/SpiderMonkey sem retreinamento completo. A questão central
investigada foi: é possível construir um único classificador que opere corretamente
em ambos os motores sem retreinamento por motor?

A resposta é afirmativa, com custo mensurável. As contribuições que sustentam essa
resposta são descritas a seguir.

\textbf{Instrumentação do SpiderMonkey com rastreamento de funções anônimas.} O
Firefox foi modificado para emitir, a cada função compilada, sua sequência de
bytecodes e metadados de identificação. O problema específico das funções anônimas
(a maioria das funções em bundles minificados) foi resolvido injetando
um identificador único como \texttt{displayName} no momento da compilação, que se
propaga para \texttt{Error.stack} e permite o join direto com as chamadas de
API capturadas pela extensão. Essa solução recuperou dados que a abordagem por
posição no código-fonte perdia em $\approx$38\% dos casos.

\textbf{Pipeline de coleta cross-engine.} Um crawler de dupla engine visita
os 10.000 sites mais populares da lista Tranco com Firefox e Chromium
instrumentados. Uma extensão de navegador compartilhada (Manifest~V3) intercepta 32
APIs de fingerprinting, normaliza propriedades exclusivas do Firefox para que
as bibliotecas de rastreamento executem o mesmo caminho de código nos dois
navegadores, e captura chamadas com resolução de frame múltiplo no call
stack. O resultado é um dataset rotulado de \textbf{68.008.403 funções JavaScript}
com 37.644 exemplos positivos de sete categorias de fingerprinting, produzido
com ambos os motores sobre o mesmo conjunto de sites.

\textbf{Representação intermediária cross-engine (contribuição principal).} Um
vocabulário de \textbf{145 tokens} (136 ativos no TF-IDF) que mapeia os opcodes nativos do V8 (198
entradas) e do SpiderMonkey (236 entradas) para uma representação compartilhada,
construído por execução pareada de 27 micro-scripts e alinhamento empírico das
sequências de bytecodes resultantes. A RI atinge \textbf{cobertura de 99,82\%} dos
tokens observados no corpus real de 68 milhões de funções, com apenas 0,18\%
mapeados para \texttt{<UNK>}. Embeddings FastText Skip-gram treinados sobre o corpus
traduzido (128 dimensões, $t=0{,}0$) apresentam coerência semântica de categoria com
gap intra/inter de +0,237.

\textbf{Avaliação comparativa de classificadores e inferência iterativa.} Cinco
famílias de modelos foram treinadas e avaliadas sobre a RI na tarefa binária
(RF, LightGBM, SVM, BiLSTM e Transformer); três dessas famílias (RF, LightGBM
e Transformer) foram também avaliadas na tarefa multilabel
(Tabela~\ref{tab:multilabel}), com a mesma divisão e as mesmas features em
todos os experimentos.
O teto de detecção baseada exclusivamente em tokens de RI (sem contexto de script)
situa-se em F1~$\approx$~0,85 (LightGBM com bigrams). O melhor
modelo sem inferência iterativa, LightGBM com bigrams e densidade de script, atinge
F1~=~0,853. A inferência iterativa, que recomputa contagens de funções
previstas em múltiplos limiares como features de contexto de script, eleva o
\textbf{resultado implantável para F1~=~0,904} sem nenhum rótulo externo. O melhor classificador multilabel, Random Forest com
\texttt{balanced\_subsample}, atinge \textbf{macro-F1~=~0,564}. Os experimentos
de transferência zero-shot confirmam que a RI é necessária mas não suficiente: o
treinamento conjunto produz um modelo com \textbf{F1~=~0,903 no Chromium e
F1~=~0,904 no Firefox} (com contexto iterativo) a partir de um único conjunto de
pesos. A diferença de
$\approx$5 pontos de F1 em relação ao ByteDefender (F1~=~0,904 vs.\
$\approx$0,950) é o custo quantificado da portabilidade entre motores, acrescido de uma tarefa estritamente mais difícil: funções anônimas incluídas,
sete categorias de fingerprinting e features exclusivamente de bytecodes
sem metadados.

As perguntas de pesquisa formuladas no Capítulo~\ref{cap:intro} são respondidas
como segue. \textbf{PQ1} (cobertura da RI $>99\%$): afirmativa, 99,82\% de
cobertura no corpus real de 68 milhões de funções. \textbf{PQ2} (detector único
cross-engine com F1 comparável ao estado da arte): afirmativa, modelo conjunto
atinge F1~=~0,903 no Chromium e F1~=~0,904 no Firefox (com contexto iterativo) sem retreinamento por motor.
\textbf{PQ3} (RI suficiente para zero-shot): negativa, transferência zero-shot
falha (F1 entre 0,149 e 0,180); treinamento conjunto é condição necessária.
\textbf{PQ4} (teto de informação): F1~$\approx$~0,85 para bytecodes puros;
F1~=~0,904 com inferência iterativa sem rótulos externos; F1~=~0,928 como limite
superior oracle.

%, , , , , , , , , , , , , , , , , , , --
\section{Limitações}
\label{sec:conclusao-limitacoes}

\textbf{Teto de informação dos bytecodes puros.} O F1~$\approx$~0,85 (F1~=~0,843
com bigrams; F1~=~0,845 com trigrams) representa o
teto para detecção baseada exclusivamente em tokens de RI sem qualquer feature de
contexto de script. O F1~=~0,853 do modelo de bigrams com densidade de script
(Tabela~\ref{tab:bin}) já incorpora um primeiro passo de contexto de
script e não é comparável ao teto de bytecodes puros. A transição de bigrams para
trigrams (+0,002~F1 com 10,7$\times$ mais features) indica que o espaço de n-grams
sobre 145 tokens está essencialmente exaurido. A inferência iterativa eleva o
resultado implantável para F1~=~0,904 incorporando contagens previstas como contexto
de script, sem consultar rótulos reais. Superar o limite superior oracle
(F1~=~0,928) requer sinais ausentes nos bytecodes: nomes de funções, contexto de
grafo de chamadas ou valores em tempo de execução.

\textbf{Rótulos raros no multilabel.} WebRTC (180 funções FP no dataset completo) e
Fonte (208) possuem suporte insuficiente para aprendizado confiável. Os melhores
resultados, F1~=~0,201 e F1~=~0,345 respectivamente, refletem a escassez de
dados, não uma limitação intrínseca da abordagem. A taxa de fingerprinting
WebRTC e de enumeração de fontes nos 10.000 sites mais populares é genuinamente
baixa.

\textbf{Cobertura de Web Workers e código lazy.} A extensão de captura não injeta
hooks em contextos de Web Worker; APIs de fingerprinting invocadas em Workers
não são registradas. Funções em pacotes de rota não ativados durante a visita do
crawler não aparecem nos logs de bytecodes.

\textbf{JavaScriptCore não instrumentado.} A RI cobre V8 e SpiderMonkey. JavaScriptCore
(Safari) não foi instrumentado; a portabilidade para um terceiro motor permanece
como hipótese de extensão, não como resultado demonstrado.

\textbf{Comparação com ByteDefender parcialmente controlada.} A composição do
dataset do ByteDefender, seleção de sites, profundidade de crawl e taxa
de positivos, não foi divulgada publicamente. O gap de $\approx$5 pontos de F1
comparando o modelo implantável (F1~=~0,904) com o ByteDefender
($\approx$0,950), inclui fatores não isoláveis além da diferença de motor,
incluindo a inclusão de funções anônimas e um espaço de rótulos mais amplo
(7 categorias vs.\ 4 do ByteDefender).

\textbf{Custo de abstração da RI não isolado.} O custo de $\approx$5 pontos de
F1 em relação ao ByteDefender é atribuído à portabilidade entre motores, mas
mistura pelo menos três fatores: (1)~perda de informação pela abstração da RI
(fusão de opcodes distintos em tokens compartilhados); (2)~ruído distribucional
do corpus conjunto (V8 e SpiderMonkey geram n-grams distintos para a mesma
construção JavaScript); e (3)~tarefa estritamente mais difícil (funções
anônimas, 7 categorias, sem metadados). Um baseline limpo para isolar o
componente~(1) seria um modelo treinado e testado \emph{no mesmo motor}, mas
sobre a RI em vez dos bytecodes brutos. Esse experimento não foi realizado; a
decomposição do custo permanece qualitativa.

%, , , , , , , , , , , , , , , , , , , --
\section{Trabalhos futuros}
\label{sec:conclusao-futuros}

\textbf{Extensão a JavaScriptCore.} A arquitetura do sistema foi projetada para
receber um terceiro motor: instrumentar o JavaScriptCore, estender o dicionário de
equivalência com os novos mapeamentos (preservando os 145 tokens existentes, 136 ativos no TF-IDF) e
retreinar o modelo conjunto. Essa extensão tornaria o sistema aplicável ao
ecossistema completo de navegadores de desktop.

\textbf{Sinais complementares para superar o teto.} A análise do
Capítulo~\ref{cap:discussao} identificou três fontes de informação que poderiam
elevar o F1 além do limite superior oracle (F1~=~0,928): nomes de funções e grafo
de chamadas, valores de propriedades em tempo de execução, e vocabulário de opcodes
mais fino (com menos fusões na RI). Uma abordagem híbrida que combine bytecodes com embeddings de nome
de função, quando disponíveis em código não minificado, é uma extensão natural.

\textbf{Coleta direcionada para rótulos raros.} WebRTC e Fonte atingiriam F1 mais
alto com suporte de treinamento adequado. Identificar e incluir sites conhecidos
por usar fingerprinting via WebRTC (ex.: serviços de videoconferência) e
enumeração de fontes (ex.: ferramentas de design web) no corpus de crawl
aumentaria o suporte positivo dessas categorias sem requerer mudanças no modelo.

\textbf{Implantação como extensão de navegador.} Os artefatos produzidos, extensão de captura de bytecodes, dicionário de RI e modelo LightGBM de 7,1~MB, são suficientes para uma prova de conceito de extensão de navegador com detecção em
tempo real. A latência de 44~µs por função é compatível com análise durante o
carregamento de página; a compilação do modelo para WebAssembly viabilizaria
execução local sem servidor.

\textbf{Análise longitudinal.} O pipeline de crawl é idempotente e
retomável. Repetir o crawl sobre os mesmos sites em intervalos regulares
permitiria rastrear a evolução das técnicas de fingerprinting ao longo do
tempo, novas APIs exploradas, mudanças em bibliotecas populares, e avaliar a
estabilidade do modelo treinado diante de concept drift.

\textbf{Web Workers.} Estender a extensão de captura para contextos de Worker (via
\texttt{ServiceWorker} e \texttt{DedicatedWorkerGlobalScope}) fecharia a lacuna de
cobertura atual e potencialmente descobriria vetores de fingerprinting não
visíveis no contexto da página principal.
